\subsection{Rollout buffers and policy versions}

The generated trajectories and their training metadata are held in a short-lived rollout
buffer. For each response, the buffer normally contains the prompt, exact sampled token
IDs, behavior-policy log-probabilities, loss masks, reward, advantage, and length. The
buffer separates generation from optimization without changing which policy generated the
data.

Three policy names have different meanings:
\begin{itemize}
  \item $\pi_{\mathrm{old}}$ is the behavior-policy snapshot that generated the current
        buffer;
  \item $\pi_\theta$ is the current trainable policy; and
  \item $\pi_{\mathrm{ref}}$ is a frozen reference policy used to penalize long-term drift.
\end{itemize}
The old policy is provenance for buffered actions. It does not mean that new rollouts
should ignore an improved current model.

In the cleanest synchronous schedule, set
$\pi_{\mathrm{old}}\leftarrow\pi_\theta$, generate a buffer, take one full-batch optimizer
step, discard the buffer, and immediately resample. The ratio is one when that loss is
evaluated, so clipping is inactive. The ratio gradient in
Equation~\ref{eq:grpo-ratio-gradient} still produces the ordinary policy gradient.

A system may instead reuse the buffer for several minibatch steps or epochs. The rationale
is computational: autoregressive rollout generation is sequential and may also wait for
tools or verifiers, whereas another training pass over stored tokens may be cheaper. After
the first optimizer step, however, $\pi_\theta$ has changed while the data still came from
$\pi_{\mathrm{old}}$. Ratios then measure the mismatch created by reuse. Reuse creates no
new information and may overfit sampled trajectories, stale their advantages, and push
many ratios into the clipped region. It is an optional compute tradeoff, not a requirement
of GRPO.
